\chapter{Retail Wholesaling}\label{ch:hf:retail-wholesaling}

A retail order to buy fifty shares is routed from a phone to a wholesaler, which fills it at a price slightly better than the best offer on any
exchange, keeps the trade, and pays the broker for the privilege. In this chapter's model, a wholesaler that gives retail orders 30\% of the
half-spread in price improvement and pays a tenth of a cent a share for them keeps 0.44 cent a share after hedging; the retail trader pays an
effective spread of 70\% of the quoted one. The same terms offered to institutional orders would lose 0.94 cent a share: the business is the
segmentation.

\section{The economics of retail flow}

Book 1, chapter 10 describes the structure: retail brokers route their customers' marketable orders to wholesalers, off-exchange market makers
that fill them as principals, usually at or inside the national best bid and offer, and pay the brokers for the flow. The wholesaler's
revenue per share has five parts:
\[
\text{net}=\underbrace{h(1-\pi)}_{\text{capture}}-\underbrace{a}_{\text{mark-out}}-\underbrace{p}_{\text{payment for flow}}
-\underbrace{c_H}_{\text{hedging}}+\underbrace{I}_{\text{inventory P\&L}},
\]
with $h$ the half-spread, $\pi$ the share of it given back as price improvement, $a$ the information in the order, $p$ the payment and $c_H$ the
cost of hedging what the wholesaler will not hold. Everything turns on $a$: retail orders are small, numerous and, over the next minute, nearly
uninformed; institutional orders are not.

\begin{definition}[Flow segmentation]\label{def:hf:retail-wholesaling:segmentation}
\emph{Flow segmentation}\index{flow segmentation} is the separation of order flow by its origin (retail, institutional, proprietary) so that each
segment trades at prices set by its own expected information and cost, rather than at a single price set by the mix; it is what allows a
wholesaler to improve on the quotes that informed flow has made wide.
\end{definition}

\begin{omfigure}
\begin{tikzpicture}[font=\footnotesize,>=stealth,box/.style={draw,rounded corners,fill=omThm!8,minimum width=2.2cm,minimum height=0.75cm,align=center}]
\node[box] (r) at (0,0) {retail trader\\(buy 50)};
\node[box] (b) at (3.2,0) {broker};
\node[box,fill=omDef!10] (w) at (6.8,0) {wholesaler};
\node[box] (x) at (11.0,1.1) {exchanges\\(NBBO)};
\node[box] (h) at (11.0,-1.1) {hedge on\\an exchange};
\draw[->] (r) -- node[above] {order} (b);
\draw[->] (b) -- node[above] {routes} (w);
\draw[->,dashed] (w.north) to[bend right=25] node[above] {payment for flow} (b.north);
\draw[->,dotted] (x.west) -- (w.north east);
\node[font=\scriptsize,anchor=south] at (8.9,0.95) {quote};
\draw[->] (w.south east) -- (h.west);
\node[font=\scriptsize,anchor=north,align=center] at (8.7,-0.95) {excess\\inventory};
\draw[->,dashed] (w.south) to[bend left=25] node[below] {fill at NBBO less improvement} (r.south);
\end{tikzpicture}

\omcaption{A retail order's path: the broker routes it to a wholesaler, which fills it as principal at the national best bid or offer improved by
part of the half-spread, pays the broker for the flow, and hedges on an exchange the inventory it will not hold. Source: Book 1, chapter 10; this
chapter.}\label{fig:hf:retail-wholesaling:path}
\end{omfigure}

\section{Segmentation and pricing}

The chapter's day (\cref{lst:hf:retail-wholesaling:flow}): 20,000 retail orders arriving about once a second, sizes lognormal around a median of 50
shares, 0.1 cent a share of information in the order's direction, sides driven by a herding sentiment that makes retail buy (or sell) in waves; the
stock quoted 2 to 3 cents wide. For comparison, a thousand institutional orders around a median of 2,000 shares, with 1.5 cents of information.
The wholesaler fills every order at the quote improved by $\pi$ of the half-spread, pays 0.1 cent a share, holds up to 10,000 shares of inventory
and hedges the excess on an exchange at the half-spread plus a 0.3-cent fee (\cref{lst:hf:retail-wholesaling:wholesale}).

\begin{omfigure}
\begin{tikzpicture}
\begin{axis}[width=11cm,height=5.4cm,xlabel={\small price improvement (\% of the half-spread)},ylabel={\small cents a share},xmin=0,xmax=80,
  ymin=-1.7,ymax=1.0,tick label style={font=\footnotesize},grid=major,grid style={gray!25},table/col sep=comma,
  legend style={font=\scriptsize,at={(0.5,-0.3)},anchor=north},legend columns=3]
\addplot[omThm,thick,mark=*,mark size=1.3pt] table[x=pi,y=retail]{figdata/hft/23-retail-wholesaling/pi.csv};
\addplot[omDef,thick,dashed,mark=square*,mark size=1.3pt] table[x=pi,y=inst]{figdata/hft/23-retail-wholesaling/pi.csv};
\addplot[gray,thick,dotted,mark=triangle*,mark size=1.3pt] table[x=pi,y=improvement]{figdata/hft/23-retail-wholesaling/pi.csv};
\draw[black,thin] (axis cs:0,0) -- (axis cs:80,0);
\legend{{wholesaler, retail flow},{wholesaler, institutional flow},{improvement to the trader}}
\end{axis}
\end{tikzpicture}

\omcaption{The wholesaler's net per share before inventory P\&L (capture less mark-out, payment for flow and hedging) on retail and on institutional
flow, and the price improvement the retail trader receives, against the share of the half-spread given back; payment for flow 0.1 cent a share,
inventory limit 10,000 shares. Data: \texttt{hf\_wholesale.by\_pi}.}\label{fig:hf:retail-wholesaling:pi}
\end{omfigure}

At no improvement the wholesaler keeps 0.85 cent a share of retail flow and loses 0.57 on institutional flow
(\cref{fig:hf:retail-wholesaling:pi}). Each 10\% of the half-spread given back costs it 0.125 cent; the retail business breaks even at 68\% and
never reaches the institutional one. That gap is the value of segmentation, and it is what pays both the improvement and the payment for flow.

\section{Price improvement and its measurement}

\begin{definition}[Internalisation rate]\label{def:hf:retail-wholesaling:rate}
The \emph{internalisation rate}\index{internalisation rate} is the share of a market maker's fills (by shares or by orders) that it keeps against
its own inventory rather than offsetting on an exchange or another venue within a short interval.
\end{definition}

Price improvement is measured against the national best bid and offer at the order's arrival. The Rule 605 statistics of Book 1, chapter 10 report it
per market centre: the effective spread (twice the distance from the mid to the fill), its ratio to the quoted spread (E/Q), and the improvement per
share. In the model at 30\% improvement, the retail trader faces a quoted spread of 2.5 cents and an effective spread of 1.75: E/Q is 0.70, and the
improvement is 0.375 cent a share.

Improvement against the quote is not the same as a good price. Schwarz, Barber, Huang, Jorion and Odean placed 85,000 simultaneous market orders
through six accounts at five brokers: the mean account-level round-trip cost ranged from 0.07\% to 0.46\%, excluding commissions, because
wholesalers gave systematically different prices for the same trades to different brokers; the variation in payment for order flow did not explain
it. The \omterm{def:hf:retail-wholesaling:rate}{internalisation rate} here is 87\%: the 13\% of the retail volume that the wholesaler hedges is the flow that arrived in waves.

\begin{dated}{2026-09}{The rules on retail execution}\label{dat:hf:retail-wholesaling:rules}
In its annual report for 2025, Virtu Financial records that in June 2025 the SEC withdrew its pending proposals for an Order Competition Rule
(proposed Rule 615 of Regulation NMS), for Regulation Best Execution and for a restriction on volume-based tiered exchange pricing, and that the
compliance date of the amendments to Rule 605 of Regulation NMS, previously around 15 December 2025, was postponed to 1 August 2026. On 3 July
2012 the SEC approved the New York Stock Exchange's Retail Liquidity Program on a pilot basis, with non-displayed retail price improvement orders
priced better than the protected best bid or offer by at least \$0.001 a share.
\end{dated}

\section{Inventory from one-sided flow}

Retail flow nets out over a day but not over minutes: when many retail traders buy the same stock at once, the wholesaler is short all of it. The
inventory limit trades hedging costs against risk (\cref{fig:hf:retail-wholesaling:limits}). With a limit of 1,000 shares the wholesaler hedges 0.47
cent a share and nets 0.20 with a daily standard deviation of 0.06; at 10,000 it hedges 0.20 and nets 0.45 with a standard deviation of 0.38; at 50,000
its hedging nearly vanishes but its daily result is dominated by inventory, with a standard deviation of 1.32 cents a share. Herding decides how much
must be hedged at all: without it, the same limit hedges 0.9\% of the volume; with the persistence of the model's sentiment, 12.9\%.

\begin{omfigure}
\begin{tikzpicture}
\begin{axis}[width=11cm,height=5.2cm,xlabel={\small inventory limit (thousand shares)},ylabel={\small net cents a share (mean $\pm$ s.d.)},
  xmode=log,log basis x=10,xtick={1,2,5,10,20,50},xticklabels={1,2,5,10,20,50},xmin=0.8,xmax=65,ymin=-1.4,ymax=1.9,
  tick label style={font=\footnotesize},grid=major,grid style={gray!25},table/col sep=comma,
  legend style={font=\scriptsize,at={(0.03,0.97)},anchor=north west}]
\addplot[omThm,thick,mark=*,mark size=1.4pt,error bars/.cd,y dir=both,y explicit] table[x=limit,y=mean,y error=sd]{figdata/hft/23-retail-wholesaling/limits.csv};
\addplot[omDef,thick,dashed,mark=square*,mark size=1.3pt] table[x=limit,y=hedge]{figdata/hft/23-retail-wholesaling/limits.csv};
\legend{{net per share, mean $\pm$ s.d.\ over 20 days},hedging cost per share}
\end{axis}
\end{tikzpicture}

\omcaption{The wholesaler's daily net per share (mean and standard deviation over twenty simulated days of retail flow) and its hedging cost per
share, against the inventory it holds before hedging on an exchange; 30\% price improvement, 0.1 cent a share paid for flow. Data:
\texttt{hf\_wholesale.limits}.}\label{fig:hf:retail-wholesaling:limits}
\end{omfigure}

\section{Rules and the debate}

The debate over payment for order flow (Book 1, chapter 10) is about who captures the value of segmentation: the wholesaler, the broker through
payments, or the retail trader through improvement. The chapter's arithmetic shows the room: 1.25 cents of half-spread, less 0.1 of information and 0.2
of hedging, leaves 0.95 cent a share to share out; the model gives 0.375 to the trader and 0.1 to the broker, and the wholesaler keeps 0.475 (0.44
after its inventory's P\&L). The withdrawn Order Competition Rule would have sent many retail orders to auctions
where other traders could compete for them; exchanges' \omterm{def:hf:retail-wholesaling:rlp}{retail liquidity programmes}, the NYSE's approved in 2012, let exchange members offer
retail orders sub-penny price improvement on the exchange itself.

\begin{definition}[Retail liquidity programme]\label{def:hf:retail-wholesaling:rlp}
A \emph{retail liquidity programme}\index{retail liquidity programme} is an exchange mechanism in which orders identified as retail can execute
against non-displayed price-improving orders that only retail orders may take, priced inside the best bid or offer by a fraction of a cent.
\end{definition}

\section{Strategy files}

\begin{strategyfile}{Equity retail internalisation}\label{strat:hf:retail-wholesaling:equity}
\sfield{Who pays you, and why} Retail traders, through the part of the spread not given back, for immediacy and a price at or inside the quote.
\sfield{Instruments and venues} US equities; brokers' routed flow; exchanges for hedging.
\sfield{Signal} The quote at arrival; the flow's segment and expected information by broker and order type.
\sfield{Sizing and execution} Fill every order at the quote improved by the agreed share; hold up to the inventory limit; hedge beyond it.
\sfield{Costs} Payment for flow, hedging (0.20 cent a share at a 10,000-share limit in the model), mark-outs, the inventory's risk.
\sfield{How it dies} Flow that is less retail than it looks, and regulation of routing; the named episode is Schwarz, Barber, Huang, Jorion and
Odean's finding that the same order earned different prices at different brokers.
\sfield{Horizon, capacity, infrastructure} Milliseconds to fill, minutes to hedge; the brokers' order flow agreements.
\sfield{Backtest honestly} The quote at arrival, the order's later mark-out by broker, herding days.
\sfield{Sources} Schwarz et al.\ (2025); Book 1, chapter 10; this chapter.
\end{strategyfile}

\begin{strategyfile}{Retail options flow through price-improvement auctions}\label{strat:hf:retail-wholesaling:options}
\sfield{Who pays you, and why} Retail options traders, whose orders carry wide spreads and little information.
\sfield{Instruments and venues} Listed options; exchange price-improvement auctions (Book 1, chapter 24), where the wholesaler's affiliate may
guarantee the order.
\sfield{Signal} The options' fair values (chapter 19) and the order's segment.
\sfield{Sizing and execution} Guarantee and improve the order in the auction; hedge delta in the stock.
\sfield{Costs} Payment for flow (higher in options than in stocks), auction allocation rules, hedging.
\sfield{How it dies} Auction rules that give competitors more of the allocation; Ernst and Spatt's evidence that option internalisation is
imperfectly competitive, which is what supports the payments.
\sfield{Horizon, capacity, infrastructure} Milliseconds; auction connectivity on every options exchange.
\sfield{Backtest honestly} Each exchange's auction allocation as of the date.
\sfield{Sources} Ernst and Spatt (2022, working paper); Book 1, chapter 24.
\end{strategyfile}

\begin{strategyfile}{Exchange retail liquidity programmes}\label{strat:hf:retail-wholesaling:rlp}
\sfield{Who pays you, and why} Retail orders that brokers send to an exchange's retail programme.
\sfield{Instruments and venues} Exchanges' retail programmes; non-displayed price-improvement orders.
\sfield{Signal} The same segmentation, on the exchange: the retail flag.
\sfield{Sizing and execution} Rest sub-penny improving orders that only retail orders can take; hedge on the lit book.
\sfield{Costs} Exchange fees; the improvement given.
\sfield{How it dies} Brokers routing retail flow to wholesalers instead; the named episode is the NYSE's programme, approved in July 2012 as a pilot.
\sfield{Horizon, capacity, infrastructure} Milliseconds; exchange membership.
\sfield{Backtest honestly} Only the retail orders that actually reached the programme.
\sfield{Sources} SEC order approving the NYSE Retail Liquidity Program (2012).
\end{strategyfile}

\section{Tutorial: fifty shares from a phone}\label{tut:hf:retail-wholesaling:tut}

\begin{tutorial}
\textbf{Goal.} Decompose a wholesaler's result on segmented flow and find how much price improvement and how much inventory it can afford.
\textbf{End state:} \cref{fig:hf:retail-wholesaling:pi,fig:hf:retail-wholesaling:limits} and the Rule 605 statistics.
\begin{tutsteps}
\item \textbf{The flow}: retail with herding, institutional with information.
\omcode{code/firm/wholesale/firm_wholesale.py}{35}{49}{Retail sizes, the herding sentiment behind their sides, and their information; institutional orders for comparison.}\label{lst:hf:retail-wholesaling:flow}
\item \textbf{The wholesaler}: capture, mark-out, payment, hedging beyond the limit, inventory P\&L.
\omcode{code/firm/wholesale/firm_wholesale.py}{60}{79}{Every order filled at the improved quote; the inventory hedged beyond its limit; the parts per share.}\label{lst:hf:retail-wholesaling:wholesale}
\item \textbf{Sweep} the improvement and the limit (\texttt{hf\_wholesale.by\_pi}, \texttt{hf\_wholesale.limits}); the Rule 605 statistics
(\texttt{firm.wholesale.rule605}).
\end{tutsteps}
\textbf{What to change next.} Let the payment differ by broker and see who gets the improvement; add price momentum to herding days; route part of
the flow to a \omterm{def:hf:retail-wholesaling:rlp}{retail liquidity programme}.
\end{tutorial}

\section{Build: the wholesaler}\label{bld:hf:retail-wholesaling:wholesale}

\begin{build}
\bfield{Purpose} Simulate segmented flow, fill it with price improvement, hedge the residual, and report the parts of the result and Rule 605-style
statistics.
\bfield{Interface} \texttt{Flow(kind, n, seed, rate, herd, info, half\_spread, spread\_jitter)}, \texttt{wholesale(flow, pi\_share, pfof, limit,
hedge\_fee, sigma\_day, seed)}, \texttt{rule605(flow, pi\_share)}.
\bfield{Rules} Cents a share; improvement as a share of the half-spread; hedging at the half-spread plus a fee.
\bfield{Acceptance tests} \texttt{code/firm/wholesale/tests/}: retail orders smaller and less informed than institutional ones; the decomposition by
hand with no noise and no limit; more improvement lowers the net; a tighter limit hedges more; E/Q equals one less the improvement share.
\bfield{Stretch} Broker-specific payments and improvement; price impact of herding; options flow and auctions.
\end{build}

\begin{omsources}
\item C.~Schwarz, B.~M.~Barber, X.~Huang, P.~Jorion, T.~Odean, The ``actual retail price'' of equity trades, \emph{Journal of Finance} 80(5), 2025,
2507--2541.
\item T.~Ernst, C.~S.~Spatt, Payment for order flow and option internalization, working paper, 2022.
\item Virtu Financial, annual report on Form 10-K for 2025.
\item US Securities and Exchange Commission, Release No.\ 34-67347 (NYSE Retail Liquidity Program), 3 July 2012.
\end{omsources}

\section{Exercises}

\begin{exercise}[$\star$]\label{exo:hf:retail-wholesaling:1}
The quote is 50.00--50.02 and the wholesaler gives 30\% of the half-spread. At what price does a retail buy order fill, and what is its effective
spread?
\end{exercise}

\begin{exercise}[$\star$]\label{exo:hf:retail-wholesaling:2}
With a half-spread of 1.25 cents, mark-out 0.1, payment 0.1 and hedging 0.2, at what improvement share does the wholesaler break even?
\end{exercise}

\begin{exercise}[$\star$]\label{exo:hf:retail-wholesaling:3}
What E/Q ratio does 30\% improvement give, and why is it independent of the spread?
\end{exercise}

\begin{exercise}[$\star\star$]\label{exo:hf:retail-wholesaling:4}
Why does herding make a wholesaler hedge fourteen times more volume at the same limit?
\end{exercise}

\begin{exercise}[$\star\star$]\label{exo:hf:retail-wholesaling:5}
Why could the wholesaler not offer institutional flow the same terms?
\end{exercise}

\begin{exercise}[$\star\star$]\label{exo:hf:retail-wholesaling:6}
Schwarz and co-authors found price dispersion across brokers not explained by payment for flow. What else could explain it?
\end{exercise}

\begin{exercise}[$\star\star\star$]\label{exo:hf:retail-wholesaling:7}
\emph{Coding.} Rerun \texttt{hf\_wholesale.limits}: which limit maximises the mean net, and which maximises the mean less the standard deviation?
\end{exercise}

\begin{exercise}[$\star\star\star$]\label{exo:hf:retail-wholesaling:8}
\emph{Find the flaw.} ``Our E/Q is 0.70, so our customers get 30\% of the spread back: the best execution in the market.''
\end{exercise}

\section{Problem: Fifty Shares from a Phone}

\begin{problem}[{Weekend problem --- fifty shares from a phone}]\label{pb:hf:retail-wholesaling:1}
A wholesaler prices its terms for a broker's retail flow.

\medskip\noindent\textbf{Part I --- The business.}
\begin{enumerate}
\item Write the wholesaler's net per share and name its parts.
\item Define \omterm{def:hf:retail-wholesaling:segmentation}{flow segmentation}.
\item Describe the retail and institutional flows of the model.
\item What does segmentation pay for?
\end{enumerate}

\medskip\noindent\textbf{Part II --- Improvement.}
\begin{enumerate}[resume]
\item Give the net per share at 0, 30\% and 70\% improvement, for retail and institutional flow.
\item Where does the retail business break even?
\item Define the \omterm{def:hf:retail-wholesaling:rate}{internalisation rate} and give the model's.
\item Give the Rule 605 statistics at 30\%.
\end{enumerate}

\medskip\noindent\textbf{Part III --- Inventory and rules.}
\begin{enumerate}[resume]
\item What does the inventory limit trade off? Give three points.
\item What does herding do to hedging?
\item Summarise the dated box.
\item Define a \omterm{def:hf:retail-wholesaling:rlp}{retail liquidity programme}.
\end{enumerate}

\medskip\noindent\textbf{Part IV --- The verdict.}
\begin{enumerate}[resume]
\item State the \emph{named result}: the wholesaler's capture per share net of payment for order flow and hedging, and the retail trader's price
improvement against the displayed spread.
\item Who captures the value of segmentation?
\item What did Schwarz and co-authors find, and what does it mean for disclosures?
\item What would an order competition rule have changed?
\item Which strategy file is most exposed to regulation?
\item How would you detect institutional flow disguised as retail?
\item What should a broker ask a wholesaler for?
\item In one sentence: what does a wholesaler buy?
\end{enumerate}
\end{problem}

\section{Interview questions}

\begin{interviewq}[$\star$ \iqroles{trader}]\label{iq:hf:retail-wholesaling:1}
Why can a wholesaler fill retail orders inside the quote when exchange market makers cannot?
\end{interviewq}

\begin{interviewq}[$\star\star$ \iqroles{researcher}]\label{iq:hf:retail-wholesaling:2}
How would you measure the information in a broker's order flow?
\end{interviewq}

\begin{interviewq}[$\star\star$ \iqroles{developer}]\label{iq:hf:retail-wholesaling:3}
Design the system that prices a retail order in 100 microseconds from the quote, the broker's terms and the inventory.
\end{interviewq}

\begin{interviewq}[$\star\star$ \iqroles{risk}]\label{iq:hf:retail-wholesaling:4}
A meme-stock day sends a wave of retail buys in one name. What limits apply, and what happens to the wholesaler's quotes?
\end{interviewq}

\begin{interviewq}[$\star\star$ \iqroles{trader}]\label{iq:hf:retail-wholesaling:5}
A broker asks for more payment for flow. How do you decide?
\end{interviewq}

\begin{interviewq}[$\star\star\star$ \iqroles{researcher}]\label{iq:hf:retail-wholesaling:6}
With herding sentiment following an AR(1) of persistence $\rho$ per order, how does the variance of the net order imbalance over $N$ orders grow
with $\rho$?
\end{interviewq}
